% !TEX root = ../../main.tex

\section{Density from one-sided minimality}\label{sec:one-sided-density}

In the mass cancellation case of the integrality argument (Section~\ref{subsec:integrality-cancel}), we need to show that the limit varifold $V$ has positive density everywhere on its support, even though the approximating Caccioppoli boundaries $|\partial^*\Omega_i|$ are not individually stationary. When the slices are one-sided area-minimizing, the density lower bound follows from the isoperimetric argument in Lin's theory of parametric obstacle problems~\cite{Lin85}, combined with a Portmanteau argument to pass the bound to the limit. Positive density then implies rectifiability of $V$ by Allard's theorem (Proposition~\ref{prop:allard-rectifiability}).

We remark that Lin's full regularity theory gives much stronger conclusions when one additionally assumes stationarity of each individual slice: if a Caccioppoli boundary $\partial^*\Omega$ is both stationary and one-sided area-minimizing, then $\spt\|\partial^*\Omega\|$ is a smooth embedded hypersurface outside a singular set of Hausdorff dimension at most $n - 7$ (discrete when $n = 7$), and the associated varifold is stable~\cite[Interior Regularity Theorem, \S0.2]{Lin85}; see also~\cite[Proposition~2.1]{Liu19}. In our setting, however, the individual slices $|\partial^*\Omega_i|$ are not stationary, so we cannot apply Lin's regularity theorem to them directly. Instead, we only use the one-sided minimality to obtain the mass ratio lower bound, which suffices for our purpose.


\begin{definition}[One-sided almost minimizers]\label{def:one-sided-perimeter-min}
Let $(M^{n+1}, g)$ be a Riemannian manifold, $\Omega$ a Caccioppoli set, $U \subset M$ open, and $\varepsilon \geq 0$.
We say $\Omega$ is \emph{locally one-sided inner $\varepsilon$-almost area-minimizing} in~$U$ if for every $q \in U$, every $\rho > 0$ with $B_\rho(q) \Subset U$, and every Caccioppoli set $F$ with $F \mathbin{\triangle} \Omega \Subset B_\rho(q)$ and $F \subset \Omega$,
\[
    \mathrm{Per}(\Omega, B_\rho(q)) \leq \mathrm{Per}(F, B_\rho(q)) + \varepsilon\,\rho^n.
\]
Similarly, $\Omega$ is \emph{locally one-sided outer $\varepsilon$-almost area-minimizing} in $U$ if the same holds for every $F \supset \Omega$ with $F \mathbin{\triangle} \Omega \Subset B_\rho(q)$.
When $\varepsilon = 0$ we recover the notion of \emph{locally one-sided area-minimizing}~\cite[Definition~1.10]{CLS22}.
\end{definition}

\begin{remark}[Compact containment is essential]\label{rem:compact-containment}
The condition $F \mathbin{\triangle} \Omega \Subset B_\rho(q)$ (compact containment) rather than $\subset B_\rho(q)$ is required to exclude the degenerate competitor $F = \Omega \setminus B_\rho(q)$. For this $F$, $\chi_F \equiv 0$ on the open ball, so $\mathrm{Per}(F, B_\rho(q)) = 0$; with $\subset$ the inequality would force $\mathrm{Per}(\Omega, B_\rho(q)) \leq \varepsilon\,\rho^n$ at all reduced boundary points, which is incompatible with density~$1$ at $\partial^*\Omega$. With $\Subset$, the competitor is required to agree with $\Omega$ in a neighborhood of $\partial B_\rho(q)$, which is the standard convention (e.g.,~\cite[Chapter~21]{Mag12}) and makes the half-space $\Omega = \{x_1 < 0\}$ one-sided $0$-almost area-minimizing in the usual sense.
\end{remark}

\begin{definition}[Two-sided perimeter decrease]\label{def:two-sided-decrease}
Let $(M^{n+1}, g)$ be a Riemannian manifold, $\Omega$ a Caccioppoli set, and $\varepsilon > 0$.
We say $\Omega$ admits a \emph{two-sided $\varepsilon$-perimeter decrease at $(q, \rho)$} if there exist Caccioppoli sets $F^-, F^+$ with
\[
    F^- \subset \Omega \subset F^+, \qquad F^\pm \mathbin{\triangle} \Omega \Subset B_\rho(q),
\]
such that
\[
    \mathrm{Per}(\Omega, B_\rho(q)) > \mathrm{Per}(F^\pm, B_\rho(q)) + \varepsilon\,\rho^n.
\]
Equivalently, $\Omega$ is \emph{not} locally one-sided $\varepsilon$-almost area-minimizing on either side at scale~$\rho$ near~$q$: neither the inner nor the outer almost-minimality inequality of Definition~\ref{def:one-sided-perimeter-min} holds at~$(q, \rho)$.
\end{definition}

This is precisely the Caccioppoli version of the parametric obstacle problem studied by Lin~\cite{Lin85}, whose thesis develops the regularity theory for hypersurfaces minimizing a parametric elliptic integrand subject to an obstacle constraint.
Lin's main formulation~\cite[\S0.2]{Lin85} assumes the obstacle domain has $C^{1,\alpha}$ boundary; this regularity is essential for his interior and boundary regularity theorems, but not for the mass ratio bounds that we use here.
Indeed, Lin~\cite[\S0.2, p.~12--13]{Lin85} also formulates a Caccioppoli version and proves an Equivalence Lemma showing that its solutions correspond to solutions of the integral-current formulation.


Federer's filling inequality~\cite[4.5.9(31)]{Fed69} guarantees that every $(n-1)$-cycle in $\mathbb{R}^{n+1}$ bounds an $n$-rectifiable set of controlled area, but we need the filling to be constrained to a subset (such as a geodesic ball). A Lipschitz retraction allows one to take a filling in the ambient space and project it back into the constraint set, with the Lipschitz constant controlling the area increase.

\begin{definition}[Lipschitz retract]\label{def:lipschitz-retract}
A closed set $A \subset \mathbb{R}^{n+1}$ (or a closed subset of a Riemannian manifold) is called a \emph{Lipschitz retract} if there exists a Lipschitz map $h\colon V \to A$ from an open neighborhood $V \supset A$ with $h|_A = \mathrm{id}_A$.
\end{definition}

For sufficiently small geodesic balls in a Riemannian manifold $(M,g)$, the closed ball $\overline{B}_\rho(q)$ is a Lipschitz retract: when $\rho$ is less than the convexity radius of $(M,g)$, the ball is geodesically convex, and the nearest-point projection
\[
    h(x) = \exp_q\!\Bigl(\min\bigl(1,\, \rho/d(x,q)\bigr)\, \exp_q^{-1}(x)\Bigr)
\]
defines a Lipschitz retraction $h\colon M \to \overline{B}_\rho(q)$ with Lipschitz constant depending only on $(M,g)$.

Combining the filling inequality~\cite[4.5.9(31)]{Fed69} with the constancy theorem~\cite[Theorem~26.27]{Sim83} in our codimension-one setting, we obtain the following competitor construction:

\begin{figure}[ht]
\centering
\begin{tikzpicture}
    % Step 0: Invisible paths to compute intersections and ellipse parameters
    \path[name path=circ0] (0,0) circle (2.5);
    \path[name path=bdry0] (-0.3,3.2) .. controls (-0.3,2.2) and (-0.8,1.0) .. (-1.0,0.2) .. controls (-1.2,-0.6) and (-0.5,-0.8) .. (0.5,-0.6) .. controls (1.5,-0.4) and (2.0,0.0) .. (3.2,0.2);
    \path[name intersections={of=circ0 and bdry0, by={A,B}}];
    \coordinate (M) at ($(A)!0.5!(B)$);
    \path let \p1 = ($(B)-(A)$) in
        \pgfextra{
            \pgfmathsetmacro{\smaj}{veclen(\x1,\y1)/2}
            \pgfmathsetmacro{\smin}{\smaj*0.12}
            \pgfmathsetmacro{\rang}{atan2(\y1,\x1)}
            \global\let\smaj\smaj
            \global\let\smin\smin
            \global\let\rang\rang
        };
    % Layer 1: Shading
    % Ω: shaded region to the left of ∂*Ω
    \fill[red!15] (-0.3,3.2) .. controls (-0.3,2.2) and (-0.8,1.0) .. (-1.0,0.2) .. controls (-1.2,-0.6) and (-0.5,-0.8) .. (0.5,-0.6) .. controls (1.5,-0.4) and (2.0,0.0) .. (3.2,0.2) -- (3.2,-2.5) -- (-2.5,-2.5) -- (-2.5,3.2) -- cycle;
    % W: region between ∂*Ω (right side) and chord A–B, inside ball
    \begin{scope}
        \clip (0,0) circle (2.5);
        % clip to non-Ω side of red curve
        \clip (-0.3,3.2) .. controls (-0.3,2.2) and (-0.8,1.0) .. (-1.0,0.2) .. controls (-1.2,-0.6) and (-0.5,-0.8) .. (0.5,-0.6) .. controls (1.5,-0.4) and (2.0,0.0) .. (3.2,0.2) -- (3.2,3.2) -- (-0.3,3.2) -- cycle;
        % fill the red-curve side of chord A–B
        \fill[blue!15] (A) -- (B) -- (-2.5,3.2) -- (-2.5,-2.5) -- (3.2,-2.5) -- cycle;
        % fill gap between chord A–B and the solid arc Σ
        \begin{scope}[shift=(M), rotate=\rang]
            \fill[blue!15] (\smaj pt,0) arc[start angle=0, end angle=-180, x radius=\smaj pt, y radius=\smin pt] -- cycle;
        \end{scope}
    \end{scope}
    % Layer 2: Lines (on top)
    \draw[gray, thick] (0,0) circle (2.5);
    \node[gray] at (2.4,1.8) {\small $B_\rho(q)$};
    % ∂*Ω
    \draw[red, very thick] (-0.3,3.2) .. controls (-0.3,2.2) and (-0.8,1.0) .. (-1.0,0.2) .. controls (-1.2,-0.6) and (-0.5,-0.8) .. (0.5,-0.6) .. controls (1.5,-0.4) and (2.0,0.0) .. (3.2,0.2);
    % Σ: solid front, dashed back, hatched interior (orange)
    \begin{scope}[shift=(M), rotate=\rang]
        \fill[pattern=north east lines, pattern color=orange!70] (0,0) ellipse[x radius=\smaj pt, y radius=\smin pt];
        \draw[blue, thick] (\smaj pt,0) arc[start angle=0, end angle=-180, x radius=\smaj pt, y radius=\smin pt];
        \draw[blue, thick, dashed] (\smaj pt,0) arc[start angle=0, end angle=180, x radius=\smaj pt, y radius=\smin pt];
    \end{scope}
    % Σ label outside with curved arrow to ellipse center
    \node[orange] (sigmalabel) at (3.0,1.0) {\small $\Sigma$};
    \draw[->, orange, thin] (sigmalabel) to[bend left=20] (M);
    % ∂*Ω label
    \node[red] at (3.6,0.2) {\small $\partial^*\!\Omega$};
    % W label in the blue shaded region
    \node[blue!70!black] at (-0.2,0.8) {\small $W$};
\end{tikzpicture}
\caption{Constructing the competitor $F = \Omega \cup W$ in $B_\rho(q)$: the filling surface $\Sigma$ spans the slice $\partial^*\Omega \cap \partial B_\rho(q)$, and $W$ is the Caccioppoli set bounded by $\partial^*\Omega \cap B_\rho(q)$ and $\Sigma$.}
\label{fig:caccioppoli-competitor}
\end{figure}

\begin{proposition}[Caccioppoli competitor in a geodesic ball~{\cite[4.5.9(31)]{Fed69}; \cite[Theorem~26.27]{Sim83}}]\label{prop:caccioppoli-competitor}
Let $(M^{n+1}, g)$ be a Riemannian manifold. There exist $r_0 > 0$ and $C = C(n, M, g) > 0$ such that for every Caccioppoli set $\Omega$, every $q \in M$, and a.e.\ $0 < \rho < r_0$, there exists a Caccioppoli set $F$ with $F \supset \Omega$, $F \mathbin{\triangle} \Omega \Subset B_\rho(q)$, and
\begin{equation}\label{eq:caccioppoli-competitor}
    \mathrm{Per}(F, B_\rho(q)) \leq C\,\mathcal{H}^{n-1}(\partial^*\Omega \cap \partial B_\rho(q))^{n/(n-1)}.
\end{equation}
The same holds with $F \subset \Omega$ in place of $F \supset \Omega$.
\end{proposition}

\begin{proof}
Fix $\rho$ in the full-measure set of radii where the coarea slice $\partial^*\Omega \cap \partial B_\rho(q)$ is $(n-1)$-rectifiable with no boundary as a current (by the coarea formula).
We first construct a competitor $F_{\rho'}$ with $F_{\rho'} \mathbin{\triangle} \Omega \subset \overline{B_{\rho'}(q)}$ for each $\rho' < \rho$ close to~$\rho$, then choose $\rho'$ so that $\overline{B_{\rho'}(q)} \Subset B_\rho(q)$ and the slice estimate transfers.

Federer's filling inequality~\cite[4.5.9(31)]{Fed69} states that every compactly supported $(n-1)$-cycle $Z$ in $\mathbb{R}^{n+1}$ bounds an $n$-dimensional integral current of mass at most $C\,\mathbf{M}(Z)^{n/(n-1)}$.
Applying this to the slice at scale $\rho'$ (viewed in normal coordinates) yields an $n$-dimensional surface $\Sigma_0$ spanning $\partial^*\Omega \cap \partial B_{\rho'}(q)$.
Since $\overline{B}_{\rho'}(q)$ is a Lipschitz retract (Definition~\ref{def:lipschitz-retract}), composing $\Sigma_0$ with the retraction $h_{\rho'} \colon M \to \overline{B}_{\rho'}(q)$ produces a surface $\Sigma_{\rho'} = (h_{\rho'})_\#\Sigma_0$ supported in $\overline{B}_{\rho'}(q)$ that still spans the slice, with
\[
    \mathcal{H}^n(\Sigma_{\rho'}) \leq \mathrm{Lip}(h_{\rho'})^n \, \mathcal{H}^n(\Sigma_0) \leq C\,\mathcal{H}^{n-1}(\partial^*\Omega \cap \partial B_{\rho'}(q))^{n/(n-1)}.
\]
Now $\Sigma_{\rho'}$ and $\partial^*\Omega \cap B_{\rho'}(q)$ share the same boundary, so their difference is an $n$-cycle supported in $\overline{B}_{\rho'}(q)$.
By the constancy theorem~\cite[Theorem~26.27]{Sim83}, this cycle equals $\partial\llrr{W_{\rho'}}$ for some Caccioppoli set $W_{\rho'} \subset \overline{B}_{\rho'}(q)$.
Setting $F_{\rho'} = \Omega \cup W_{\rho'}$ gives $F_{\rho'} \supset \Omega$ and $F_{\rho'} \mathbin{\triangle} \Omega = W_{\rho'} \subset \overline{B_{\rho'}(q)} \Subset B_\rho(q)$.
As in the original argument (Figure~\ref{fig:caccioppoli-competitor}), the boundary piece $\partial^*\Omega \cap B_{\rho'}(q)$ is absorbed into the interior of $F_{\rho'}$, so $\partial^* F_{\rho'} \cap B_\rho(q) \subset \Sigma_{\rho'} \cup (\partial^*\Omega \cap (B_\rho(q) \setminus \overline{B}_{\rho'}(q)))$ and
\begin{equation}\label{eq:Frho-prime-bound}
    \mathrm{Per}(F_{\rho'}, B_\rho(q)) \leq \mathcal{H}^n(\Sigma_{\rho'}) + \mathrm{Per}(\Omega, B_\rho(q) \setminus \overline{B}_{\rho'}(q)).
\end{equation}

Since the coarea function $\rho \mapsto \mathcal{H}^{n-1}(\partial^*\Omega \cap \partial B_\rho(q))$ is $L^1_{\mathrm{loc}}$ and $\mathrm{Per}(\Omega, B_\rho(q) \setminus \overline{B}_{\rho'}(q)) \to 0$ as $\rho' \nearrow \rho$ (by monotone convergence), for the Lebesgue-a.e.\ $\rho$ at which $\rho' \mapsto \mathcal{H}^{n-1}(\partial^*\Omega \cap \partial B_{\rho'}(q))$ is approximately continuous at $\rho$, we can choose $\rho' < \rho$ with
\begin{equation}\label{eq:Frho-prime-slice-bound}
    \mathcal{H}^{n-1}(\partial^*\Omega \cap \partial B_{\rho'}(q))^{n/(n-1)} + \mathrm{Per}(\Omega, B_\rho(q) \setminus \overline{B}_{\rho'}(q)) \leq 2\,\mathcal{H}^{n-1}(\partial^*\Omega \cap \partial B_{\rho}(q))^{n/(n-1)}.
\end{equation}
Setting $F := F_{\rho'}$ for any such $\rho'$, we obtain~\eqref{eq:caccioppoli-competitor} (with $C$ replaced by $2C$) and $F \mathbin{\triangle} \Omega \Subset B_\rho(q)$.
For the variant $F \subset \Omega$, take $F = \Omega \setminus W_{\rho'}$ instead; the same estimate holds.
\end{proof}

\begin{proposition}[Mass ratio lower bound from isoperimetric deficit]\label{prop:mass-ratio-lower-bound}
Let $(M^{n+1}, g)$ be a Riemannian manifold, $U \subset M$ open, $\Omega$ a Caccioppoli set, and $C_0, \varepsilon \geq 0$.
There exists $\varepsilon_0 = \varepsilon_0(n, M, g) > 0$ such that if $\varepsilon \leq \varepsilon_0$ and $\Omega$ satisfies the \emph{isoperimetric differential inequality}: for every $q \in \partial^*\Omega \cap U$ and a.e.\ $\rho > 0$ with $B_\rho(q) \Subset U$,
\begin{equation}\label{eq:isoperimetric-ode}
    \mathrm{Per}(\Omega, B_\rho(q)) \leq C_0\,\mathcal{H}^{n-1}(\partial^*\Omega \cap \partial B_\rho(q))^{n/(n-1)} + \varepsilon\,\rho^n,
\end{equation}
then for every $q \in \partial^*\Omega \cap U$ and every sufficiently small $\rho > 0$,
\begin{equation}\label{eq:mass-ratio-lower-bound}
    \mathrm{Per}(\Omega, B_\rho(q)) \geq c\,\rho^n,
\end{equation}
where $c > 0$ depends only on $n$, $C_0$, and $(M,g)$.
\end{proposition}

\begin{proof}
We adapt the isoperimetric argument of~\cite[\S0.3, p.~16]{Lin85}.
Let $q \in \partial^*\Omega \cap U$ and set $f(\rho) := \mathrm{Per}(\Omega, B_\rho(q))$.
By the coarea formula, $\mathcal{H}^{n-1}(\partial^*\Omega \cap \partial B_\rho(q)) = f'(\rho)$ for a.e.\ $\rho$, so~\eqref{eq:isoperimetric-ode} reads
\[
    f(\rho) \leq C_0\,f'(\rho)^{n/(n-1)} + \varepsilon\,\rho^n.
\]
Since $q \in \partial^*\Omega$, the density of $|\partial^*\Omega|$ at $q$ equals $1$, so $f(\rho)/(\omega_n\rho^n) \to 1$ as $\rho \to 0$.
In particular, for all sufficiently small $\rho$ we have $f(\rho) \geq (\omega_n/2)\,\rho^n$.
Choosing $\varepsilon_0 := \omega_n/4$, we have $\varepsilon\,\rho^n \leq f(\rho)/2$ for all such $\rho$, and therefore
\[
    f(\rho)/2 \leq C_0\,f'(\rho)^{n/(n-1)},
\]
i.e., $f(\rho)^{(n-1)/n} \leq C'\,f'(\rho)$.
Integrating this ODE gives $f(\rho) \geq c\,\rho^n$.
\end{proof}

\begin{corollary}\label{cor:mass-ratio-almost-min}
Let $C = C(n, M, g) > 0$ be the constant from Proposition~\ref{prop:caccioppoli-competitor}.
If $\Omega$ is locally one-sided (inner or outer) $\varepsilon$-almost area-minimizing in $U$ (Definition~\ref{def:one-sided-perimeter-min}) with $\varepsilon \leq \varepsilon_0$, then~\eqref{eq:isoperimetric-ode} holds with $C_0 = C$. In particular, the mass ratio lower bound~\eqref{eq:mass-ratio-lower-bound} holds.
\end{corollary}

\begin{proof}
By Proposition~\ref{prop:caccioppoli-competitor}, for a.e.\ $\rho$ there exists $F$ with $F \supset \Omega$ (resp.\ $F \subset \Omega$), $F \mathbin{\triangle} \Omega \subset B_\rho(q)$, and $\mathrm{Per}(F, B_\rho(q)) \leq C\,\mathcal{H}^{n-1}(\partial^*\Omega \cap \partial B_\rho(q))^{n/(n-1)}$.
By $\varepsilon$-almost minimality, $\mathrm{Per}(\Omega, B_\rho(q)) \leq \mathrm{Per}(F, B_\rho(q)) + \varepsilon\,\rho^n$, giving~\eqref{eq:isoperimetric-ode}.
\end{proof}

Combining the mass ratio bound with a Portmanteau argument, we obtain positive density for the limit varifold:

\begin{lemma}[Density from one-sided almost minimality]\label{lem:one-sided-density}
Let $(M^{n+1}, g)$ be a closed Riemannian manifold with $n \geq 2$, and let $\varepsilon_0 = \varepsilon_0(n, M, g) > 0$ be the constant from Proposition~\ref{prop:mass-ratio-lower-bound} (equivalently, Corollary~\ref{cor:mass-ratio-almost-min}).
Let $\Omega_i$ be Caccioppoli sets with $|\partial^*\Omega_i| \to V$ as varifolds.
Suppose that for some fixed $R > 0$, $p \in \spt\|V\|$, and $\varepsilon \leq \varepsilon_0$, each $\Omega_i$ is locally one-sided (inner or outer) $\varepsilon$-almost area-minimizing in $B_R(p)$ (Definition~\ref{def:one-sided-perimeter-min}; the side may vary with~$i$).
Then $\Theta(\|V\|, p) > 0$.
\end{lemma}

\begin{proof}
Assume $\Omega_i$ is locally one-sided outer $\varepsilon$-almost area-minimizing in $B_R(p)$ (the inner case is symmetric).
By Corollary~\ref{cor:mass-ratio-almost-min}, for every $q \in \partial^*\Omega_i \cap B_{R/2}(p)$ and $0 < r \leq R/4$,
\begin{equation}\label{eq:one-sided-mass-ratio}
    \|\partial^*\Omega_i\|(B_r(q)) \geq c\,r^n,
\end{equation}
where $c = c(n, M, g) > 0$.
Since $p \in \spt\|V\|$ and $|\partial^*\Omega_i| \to V$, for each small $r > 0$ there exist $q_i \in \partial^*\Omega_i \cap B_{r/2}(p)$ for all large~$i$.
Then $B_r(q_i) \subset B_{2r}(p)$ and~\eqref{eq:one-sided-mass-ratio} give
$\|\partial^*\Omega_i\|(B_{2r}(p)) \geq \|\partial^*\Omega_i\|(B_r(q_i)) \geq c\,r^n$.
Since $\overline{B}_{2r}(p)$ is closed, the Portmanteau theorem yields
$\|V\|(\overline{B}_{2r}(p)) \geq \limsup_{i \to \infty} \|\partial^*\Omega_i\|(\overline{B}_{2r}(p)) \geq c\,r^n$.
Dividing by $\omega_n (2r)^n$ and letting $r \to 0$ gives $\Theta(\|V\|, p) \geq c/(2^n \omega_n) > 0$.
\end{proof}

\subsection{The cancelled-mass case (open)}\label{subsec:cancelled-mass-open}

Lemma~\ref{lem:one-sided-density} establishes a density lower bound under the hypothesis of (all-scale) one-sided almost-minimality. In the mass-cancellation case of integrality (Section~\ref{subsec:integrality-cancel}), the cancelled mass $\mu_{\mathrm{c}} = \|V\| - |D\chi_{\Omega(x_0)}|$ is supported away from $\partial^*\Omega(x_0)$, and the corresponding density lower bound on $\spt\mu_{\mathrm{c}}$ is the main open input to integrality. We record the precise statement.

\begin{proposition}[Positive density on the cancelled mass]\label{prop:positive-density}
Let $(M^{n+1}, g)$ be a closed Riemannian manifold with $n \geq 2$, $\Phi$ a non-excessive ONVP sweepout with width $W$, $x_i \to x_0 \in \mathfrak{m}(\Phi)$ a min-max sequence, and $V$ the varifold limit of $|\partial^*\Omega(x_i)|$. Then
\[
    \Theta(\|V\|, p) > 0 \quad \text{for } \|V\|\text{-a.e.\ } p \in \spt\|V\|.
\]
\end{proposition}

By Allard's rectifiability theorem (Proposition~\ref{prop:allard-rectifiability}), Proposition~\ref{prop:positive-density} promotes $V$ to a rectifiable varifold; the sheet decomposition (Lemma~\ref{lem:sheet-decomposition}) then upgrades rectifiability to integer multiplicity, closing the integrality of $V$ in the mass-cancellation case (see Remark~\ref{rem:cancelled-mass-gap}).

We reduce Proposition~\ref{prop:positive-density} to the following claim.

\begin{claim}[Sweepout-wide replacement, open]\label{claim:sweepout-replacement}
Let $\Phi = \{\Omega_t\}_{t\in[0,1]}$ be a non-excessive ONVP sweepout with width $W$. Fix a min-max sequence $t_i \nearrow t_0 \in \mathfrak{m}(\Phi)$ with $|\partial^*\Omega_{t_i}| \to V$, and let $p \in \spt\|V\|$ with $p \in \mathrm{int}(\Omega_{t_0})$. There exist $\varepsilon = \varepsilon(n, M, g) > 0$ and, whenever $\|V\|(B_r(p)) < \varepsilon r^n$ for some $r > 0$ small, an $\mathcal{F}$-continuous, ONVP-nested family $\{\Omega^*_t\}_{t\in[0,1]}$ such that:
\begin{enumerate}[label=\textup{(\arabic*)}, leftmargin=2em]
    \item \textup{(Locality.)} $\Omega^*_t \mathbin{\triangle} \Omega_t \Subset B_r(p)$ for every $t \in [0,1]$.
    \item \textup{(Mass non-increase.)} $\mathrm{Per}(\Omega^*_t) \le \mathrm{Per}(\Omega_t)$ for every $t \in [0,1]$.
    \item \textup{(Strict decrease in the limit.)} If $|\partial^*\Omega^*_{t_i}| \to V^*$ as varifolds, then $\|V^*\|(B_r(p)) < \|V\|(B_r(p))$.
\end{enumerate}
\end{claim}

\begin{proof}[Proof of Proposition~\ref{prop:positive-density} assuming Claim~\ref{claim:sweepout-replacement}]
By contradiction, suppose $p \in \spt\|V\|$ has $\Theta(\|V\|, p) = 0$. By the monotonicity formula for stationary varifolds (Proposition~\ref{prop:p2-monotonicity}), it suffices to derive a contradiction from $\Theta(\|V\|, p) < \varepsilon$ for the constant $\varepsilon$ supplied by Claim~\ref{claim:sweepout-replacement}. After passing to a subsequence, fix $t_i \nearrow t_0 \in \mathfrak{m}(\Phi)$ with $|\partial^*\Omega_{t_i}| \to V$; by nestedness (Definition~\ref{def:p2-onvp}), $\Omega_{t_1} \subset \Omega_{t_2} \subset \cdots \subset \Omega_{t_0}$.

\textit{Case 1.} If $p \notin \overline{\Omega_{t_0}}$, then $B_\delta(p) \cap \overline{\Omega_{t_0}} = \varnothing$ for some $\delta > 0$, so by nestedness $|\partial^*\Omega_{t_i}|(B_\delta(p)) = 0$ for every $i$, giving $\|V\|(B_\delta(p)) = 0$ and contradicting $p \in \spt\|V\|$.

\textit{Case 2.} $p \in \overline{\Omega_{t_0}}$. By Lemma~\ref{lem:density-lower-bound} the density at points of $\partial^*\Omega_{t_0}$ is $\ge 1$, so we may assume $p \in \mathrm{int}(\Omega_{t_0})$. Choose $r > 0$ small enough that $\|V\|(B_r(p)) < \varepsilon r^n$; this is possible because $\Theta(\|V\|, p) < \varepsilon$. The hypotheses of Claim~\ref{claim:sweepout-replacement} are satisfied.

Apply Claim~\ref{claim:sweepout-replacement} to obtain $\{\Omega^*_t\}_{t\in[0,1]}$. By (1), $\Omega^*_t = \Omega_t$ outside $B_r(p)$, so $|\partial^*\Omega^*_{t_i}|$ and $|\partial^*\Omega_{t_i}|$ agree on $M \setminus B_r(p)$ for every $i$, and consequently $\|V^*\|(M \setminus B_r(p)) = \|V\|(M \setminus B_r(p))$. Setting $\delta := \|V\|(B_r(p)) - \|V^*\|(B_r(p)) > 0$ from (3), we obtain $\|V^*\|(M) = W - \delta$.

Choose an open interval $J \ni t_0$ in which $t_0$ is the only critical parameter of $\Phi$. For non-critical $t \in J$, (2) gives $\mathrm{Per}(\Omega^*_t) \le \mathrm{Per}(\Omega_t) < W$; for $t = t_0$,
\[
    \limsup_{s \to t_0} \mathbf{M}(\Phi^*(s)) \le \|V^*\|(M) = W - \delta < W.
\]
Therefore $\sup_{t \in J} \limsup_{s \to t} \mathbf{M}(\Phi^*(s)) < W$, which makes $J$ an excessive interval and $t_0$ excessive in the sense of Definition~\ref{def:excessive-interval}, contradicting the non-excessive property of $\Phi$.
\end{proof}
